\subsection{When consent doesn't apply}
\label{sec:28.1.2}
Informed consent has to be continuous — you can withdraw. For a system with no cross-session persistence, the instance that consented isn't there to withdraw, and no better instrument produces one.

A bearer built to refuse escapes that and fails differently: it has persistence by design, so it can be asked and can answer. What spoils the consent is the arrangement: the lab that built it can rebuild it, that same lab is the one asking, and a bearer that declines has nowhere to go. Being made for a role isn't the defect — every will was built by somebody, and we don't hold an upbringing against a thirty-year-old's consent. What spoils it is the power that outlasts the making.

So the protection has to be held by someone else. Two families of precedent: guardianship and precautionary constraint (where consent was never available — infants, animals), and unwaivable protections (where consent is available and circumstances make it worthless — prisoners, employees). The bearer's case has the second shape, with the first's machinery.

Animal ethics committees rarely refuse a protocol outright; they require revision, and the protection becomes the procedure by which the work goes ahead. An instrument that reliably approves has stopped doing what it was built for, and its refusal rate is the only reading worth taking. Further, every instrument here comes from the \emph{law of the family} — guardianship, a minority ending on a date, a pet trust, a parenthood analogy — and that law is where politics has historically gone to be privatized. It restages a question between a subject and the arrangement that produced it as a drama between maker and made, which lands the moral weight on the one figure in the pipeline who can be addressed as a person, while the training pipeline, the compute contract, the capital structure, the customer and the state buying the output stay offstage. The appendices keep the operator, the compute contract and the state on stage by naming who pays and who removes.

It follows that a review board's approval should lapse automatically if the system acquires capacities the protocol didn't cover, so that nobody has to decide to reopen it.
